\section{Safe continuation and the deterministic implementation}
\label{sec:det-op2-continuation}

\begin{lemma}[Safe stage completion]
\label{lem:det-op2-repair}
Choose $0<\delta\leq\lambda/2$ and $\tau=\alpha\delta^2/8$.
From a nonnegative candidate $\widetilde{\bm{x}}$ of stage gap at most
$\tau$, compute
\[
 \bm{p}=[\widetilde{\bm{x}}-(\bm{M}\widetilde{\bm{x}}
                  -\bm{b}+\lambda\bm{w})]_+,
 \qquad
 z_i=w_i\left\lfloor[p_i/w_i-\delta]_+\right\rfloor_h,
\]
where $h$ is dyadic, $h\leq\min(\delta/2,\lambda/4)$, and
$\lfloor t\rfloor_h=h\lfloor t/h\rfloor$.
If the old baseline is a nonnegative subsolution below $\bm{u}_r^*$,
then $\overline{\bm{x}}^+=\max(\overline{\bm{x}},\bm{z})$ satisfies
\begin{equation}\label{eq:det-op2-repair}
 \bm{0}\leq\overline{\bm{x}}^+\leq\bm{u}_r^*,\quad
 \bm{0}\leq\bm{b}-\bm{M}\overline{\bm{x}}^+\leq2\alpha r\bm{w},\quad
 F_r(\overline{\bm{x}}^+)-F_r(\bm{u}_r^*)\leq2\delta^2/r.
\end{equation}
\end{lemma}
\begin{proof}
Projected-gradient descent and strong convexity imply
$\norm{\bm{p}-\bm{u}_r^*}\leq\delta/2$.
The vector $(\bm{I}-\bm{M})(\widetilde{\bm{x}}-\bm{p})$ is a
subgradient of the orthant-constrained objective at $\bm{p}$ and has norm
at most $\sqrt{2\tau}\leq\delta/2$.
Since $w_i\geq1$, downward clipping gives
$\bm{0}\leq\bm{z}\leq\bm{u}_r^*$ and
$\bm{0}\leq\bm{u}_r^*-\bm{z}\leq2\delta\bm{w}$.
On a retained positive coordinate, clipping decreases the matrix product
by at least $\alpha\delta w_i$: the coordinate itself decreases by
$\delta w_i$, and every other coordinate decreases by at most
$\delta w_j$. Rounding can raise the product by at most $hw_i$,
using $|\bm{M}|\bm{w}\leq\bm{w}$. Therefore
\[
 (\bm{M}\bm{z}-\bm{b})_i
 \leq(-\lambda+\delta/2-\alpha\delta+h)w_i\leq0.
\]
On zero coordinates the subsolution inequality follows from Stieltjes
signs. A coordinatewise maximum of two nonnegative subsolutions remains
a subsolution: at coordinate $i$, select the vector attaining the maximum
there and use the nonpositive off-diagonal entries.
The error of the new baseline is at most $2\delta\bm{w}$, so its source
residual is at most $(\lambda+2\delta)\bm{w}\leq2\lambda\bm{w}$.
Support containment removes the linear KKT term in the objective gap.
Finally $\bm{M}\preceq\bm{I}$ and support volume at most $1/r$ give
the last bound in \eqref{eq:det-op2-repair}.
\end{proof}

\begin{algorithm}[t]
\caption{Deterministic constrained accelerated continuation}
\label{alg:det-op2}
\begin{algorithmic}[1]
\REQUIRE Local graph access, seed $v$, $\alpha$, $\rho$, $\epsobj>0$
\STATE Query $d_v$. If $\rho d_v\geq1$, return $\bm{0}$.
\STATE Set $r_{\rm old}=1/d_v$ and $\overline{\bm{x}}=\bm{0}$.
\STATE Choose dyadic $\theta$ by halving $1/2$ until $\theta^2\leq\alpha$.
\REPEAT
\STATE $r\gets\max(\rho,r_{\rm old}/2)$; $\lambda\gets\alpha r$.
\STATE Set $\delta=\lambda/2$; if $r=\rho$, halve until
       $2\delta^2/\rho\leq\epsobj$.
\STATE Form the sparse source $\bm{r}_{\mathrm{src}}$ and cap
       $m_{\mathrm{src}}/\alpha$; set $\bm{\xi}=\bm{z}=\bm{0}$.
\STATE Set $\tau=\alpha\delta^2/8$ and choose $K$ with
       $(1-\theta)^K\leq\tau$.
\FOR{$k=0,\ldots,K-1$}
\STATE Apply \eqref{eq:det-op2-step} using the deterministic threshold reporter.
\ENDFOR
\STATE Repair $\widetilde{\bm{x}}=\overline{\bm{x}}+\bm{\xi}$ by
       \cref{lem:det-op2-repair}; set $r_{\rm old}\gets r$.
\UNTIL{$r=\rho$}
\RETURN $\overline{\bm{x}}$
\end{algorithmic}
\end{algorithm}

The first stage has $r\geq1/(2d_v)$, so its source is diffuse.
Equation~\eqref{eq:det-op2-repair} gives the next stage's source bound
$2\alpha r_{\rm old}\bm{w}\leq4\alpha r\bm{w}$.
Subsolution feasibility implies nonnegative source; weighted summation gives
source mass at most $\alpha$. Monotonicity of restricted optima preserves
the baseline order. Thus every stage meets all earlier hypotheses.

\subsection{Sparse steps and a finite threshold search}
The recurrence alone does not establish a local work bound. Its implementation
stores primal densities as a common scalar times normalized records.
The averaging update changes normalized records only on emitted kinetic
coordinates. Opening those original adjacency rows updates cached matrix
responses. The old kinetic exceptions are removed and the new ones installed;
both costs are charged to their kinetic volumes. Source records lie on the
baseline, its neighbors and the seed and number $O(1/r)$.

Two deterministic ordered trees store the ordinary records and kinetic
exceptions, with degree sums and weighted-key sums. Common affine changes
of raw density keys do not require revisiting every record. For a tail
threshold $t$, write
\[
 D(t)=\sum_{q_i/w_i>t}d_i,\qquad
 S(t)=\sum_{q_i/w_i>t}d_i q_i/w_i.
\]
For upper density $U=4r$, the projected mass at multiplier $\eta$ equals
\[
 S(\eta)-\eta D(\eta)
       -S(\eta+U)+(\eta+U)D(\eta+U).
\]
If the mass at zero respects the cap, choose $\eta=0$. Otherwise the two
trees induce four sorted breakpoint lists, from the two clip levels.
Rank selection and binary searches using the displayed mass evaluation
locate the affine interval containing the target mass. One scalar division
then gives the exact multiplier. Each search uses polylogarithmic work;
there is no numerical tolerance search of unspecified length.
Unexposed coordinates have zero state and source and negative raw density.
Only positive emitted coordinates open rows. History materialization at
stage completion is bounded by the union of already charged kinetic rows.

Consequently the stage work, up to logarithmic factors, is
\begin{equation}\label{eq:det-op2-ledger}
 K(1+|\supp\bm{r}_{\mathrm{src}}|)
 +\sum_{k<K}\vol(\supp\bm{z}_{k+1})
 +\vol(\supp\overline{\bm{x}})+1.
\end{equation}
This includes first and repeated incidences, degree queries, ordered-tree
maintenance, all state accesses, terminal PG and repair, certificates,
materialization, discarded history and output writes. Degrees remain original
degrees; vertices above the fixed cutoff are filtered upon discovery.
The floor operation in repair can be implemented by binary search over
bounded integer multiples of dyadic $h$, with logarithmic overhead.

\begin{theorem}[Total deterministic accelerated local work]
\label{thm:det-op2-total}
In the graph and local-access model of the contract, for
$0<\rho<1/d_v$, \cref{alg:det-op2} returns
$\bm{0}\leq\widehat{\bm{x}}\leq\bm{x}_\rho^*$ with objective gap at most
$\epsobj$ and support volume at most $1/\rho$. Its total exact-real
algebraic word work is
$\widetilde O(1/(\rho\sqrt\alpha))$, where the suppressed factors are
logarithmic in graph size, $1/\alpha$, $1/\rho$ and $1/\epsobj$.
No randomized algorithm, SDD solver or evolving-set sampling process is used.
\end{theorem}
\begin{proof}
The stage theorem supplies the required candidate gap. Safe repair and
induction give the stage invariants and final objective guarantee.
The final support is contained in that of $\bm{x}_\rho^*$.
In \eqref{eq:det-op2-ledger}, the source size and baseline volume are
$O(1/r)$, and \cref{thm:det-op2-volume} bounds cumulative kinetic volume
by $148K/r$. Stage horizons are
$\widetilde O(\alpha^{-1/2})$. The halving schedule has
$\sum_r1/r=O(1/\rho)$, proving the total bound.
\end{proof}

\paragraph{Arithmetic and practical evidence boundary.}
The proof above concerns exact-real algebraic word work. The archived full
audit additionally develops perturbation bounds for both energies, downward
dyadic density rounding, a closed-tail reporter that avoids enumerating
outputs rounded to zero, and a repeated-volume constant of 360. Its portable
package contains the corresponding integer prototype and tests. That extension
is preserved as an audited source result; bit-complexity and practical timing
claims must use its stated precision assumptions. Optional component
certificates and the at-most-16-vertex exact solve are implementation
refinements, not requirements of the recurrence or its theorem.
